\par\noindent\textbf{(1)} 球面度量 \(ds^2=R^2(d\theta^2+\sin^2\theta d\phi^2)\)。沿赤道 \(\theta=\pi/2\) 取 \(\phi=s/R\)，则 \(g_{\phi\phi}(d\phi/ds)^2=R^2/R^2=1\)。其纬度方程中 \(\sin\theta\cos\theta=0\)，满足测地线条件。
\par\noindent\textbf{(2)} 球面Gaussian曲率常数为 \(K_G=1/R^2\)，二维横向Jacobi方程 \(J''+K_GJ=0\)，通解 \(J(s)=A\cos(s/R)+B\sin(s/R)\)。该公式量化相邻自由轨道的会聚发散，与球面正曲率直接关联。
\par\noindent\textbf{(3)} 若在起点相交 \(J(0)=0\)，则 \(A=0\)。对非平凡邻轨 \(B\ne0\)，下一零点位于 \(s=\pi R\)。这是球面对径的共轭点；多条大圆在两对径点相遇，并非所有短程轨道均唯一。
\par\noindent\textbf{(4)} 当 \(0<s<\pi R\) 时大圆短弧是两点间最短路径；超过首个共轭点后仍可沿同一大圆保持测地线方程，却不再全局最短。共轭点体现变分二阶型出现零方向，不能由单纯Euler–Lagrange驻值性推断全局极小。
